\subsection*{流动场的一根箭头还不够}
向量场在每个位置给出速度。要判断附近区域膨胀还是旋转，不能只看中心那根箭头，还要看周围箭头怎样变化。放入一个很小的方框，比较各边的流入与流出，就能测量局部净扩张趋势。

对 $F=(P,Q)$，左右边的水平流量差主要来自 $P_x$，上下边的垂直差来自 $Q_y$。除以小面积再取极限，得到散度 $P_x+Q_y$。它关心的是速度变化带来的净流出，而不是速度本身大不大。

\subsection*{环流与膨胀是不同测量}
再沿小方框边界绕一圈，计算速度沿切向的贡献。相对边的差产生 $Q_x-P_y$，这是平面旋度。场 $(-y,x)$ 让小桨轮转动，却不膨胀，所以散度零、旋度非零。

场 $(2x,-2y)$ 则横向拉伸、纵向压缩，瞬时面积效应抵消，散度为零，也没有局部环流。零散度不是“什么都没发生”，而是两方向的面积效应恰好平衡。对这种线性场，追踪小方形会比看单根箭头清楚。

这些量都是局部一阶信息。把它们与整个区域的边界流量联系起来，需要散度定理或斯托克斯定理及其条件。这里只建立局部测量的意义；遇到全局势函数问题，还必须检查定义区域，而不能仅凭某一点或某个局部公式作结论。
